\chapter{Methodology}
\label{ch:methodology}

This chapter presents the engineering approach behind the Multi-Sport Club
Management System (MSCMS), commercially branded \emph{Sportify}. It documents the
software-development methodology, the elicited functional and non-functional
requirements, the cloud-native microservices architecture, the decomposition of the
back-end into domain services, the security and authentication design, the data
model, the domain (class) model, the technology stack, an overview of the
representative API surface, and the principal behavioural flows of the platform.
Whereas earlier chapters motivated the problem and surveyed comparable systems, the
purpose of the present chapter is to explain \emph{how} the platform was designed and
constructed so that the artefacts reported in Chapter~\ref{ch:results} can be
understood as the deliberate outcome of these decisions rather than incidental
features.

The chapter is organised to move from process to product and from the abstract to
the concrete. Section~\ref{sec:methodology-process} describes the development
process and the engineering practices that supported it.
Section~\ref{sec:requirements} sets out the functional and non-functional
requirements that defined the target. Sections~\ref{sec:architecture}
and~\ref{sec:decomposition} describe the macro-architecture and the way the back-end
was decomposed into autonomous services. Section~\ref{sec:security} treats security
as a first-class architectural concern. Sections~\ref{sec:data-design}
and~\ref{sec:domain-model} describe the persistent data model and the in-memory
domain model that it materialises. Section~\ref{sec:api-surface} surveys the
representative HTTP surface exposed by each service, and
Section~\ref{sec:tech-stack} records the technology choices. Finally,
Section~\ref{sec:flows} develops the system's principal behavioural flows in detail,
expressing each as a numbered sequence of interactions between the participating
components.

\section{Development Methodology}
\label{sec:methodology-process}

MSCMS was developed following an \emph{agile, iterative and incremental} process
model~\cite{agile}. A purely sequential (waterfall) plan was considered unsuitable
because the project's scope deliberately evolved over time: the system began as a
conventional create--read--update--delete (CRUD) club-administration tool and matured
into a SofaScore/ESPN-style sports platform with live-match coverage, real
competition data and real-time team communication. An agile model accommodated this
trajectory, allowing requirements to be reprioritised as the product vision sharpened
and as feedback was gathered after each working release.

The work was organised into short, time-boxed iterations, each producing a
demonstrable vertical slice that cut through every architectural layer---from the
PostgreSQL schema and the Spring Boot service, through the API Gateway, to the
Next.js user interface. This vertical-slice discipline ensured that every iteration
ended with an integrated, runnable system rather than a set of disconnected
components. The principal virtue of the approach for a project of this size was
\emph{de-risking}: each iteration validated an end-to-end path---authentication
through the gateway, persistence in the correct database, propagation of a
cross-domain event, and rendering in a role-appropriate view---so that integration
defects surfaced early, while a slice was small enough to reason about, rather than
late, when many subsystems had already been assumed to work. Representative
iterations are summarised in Table~\ref{tab:iterations}.

\begin{table}[htbp]
\centering
\caption{Representative development iterations and their primary deliverables.}
\label{tab:iterations}
\renewcommand{\arraystretch}{1.3}
\begin{tabular}{|p{0.05\textwidth}|p{0.30\textwidth}|p{0.52\textwidth}|}
\hline
\textbf{It.} & \textbf{Theme} & \textbf{Primary deliverable} \\
\hline
1 & Identity \& foundation & Keycloak realm, JWT login, API Gateway, Eureka and Config Server, the user-management service. \\
\hline
2 & Squad \& people & Player-management service: sports, teams, rosters, contracts, transfers, scouting prospects. \\
\hline
3 & Training \& matches & Training-match service: sessions, plans, drills, attendance, match fixtures, formations and line-ups. \\
\hline
4 & Competitions overhaul & Database-backed La Liga and the new-format UEFA Champions League (league phase plus knockout bracket). \\
\hline
5 & Live-match engine & Broadcast-style phase machine, extra time and penalties for knockout ties, in-match injury events. \\
\hline
6 & Medical workflow & Treatment Room and the staged per-player injury journey from injury to recovered. \\
\hline
7 & Communication & Real-time team group chat over WebSocket, event-driven notifications and the alert centre. \\
\hline
8 & Governance \& polish & Per-role RBAC, persistent settings with Keycloak-verified password change, and the Sportify rebrand. \\
\hline
\end{tabular}
\end{table}

Several engineering practices reinforced the iterative process. Version control was
maintained throughout, with the back-end, front-end and machine-learning code bases
kept under a single repository to simplify cross-cutting changes. A
container-first philosophy meant that every service was runnable through Docker
Compose from the earliest iterations, removing the ``works on my machine'' class of
defects and giving every iteration a reproducible integration environment. Finally,
single-source-of-truth artefacts---most notably the front-end permission map and the
back-end seed-identity constants---were introduced so that cross-cutting concerns
could evolve consistently without scattering duplicated logic across modules.

Two further practices deserve emphasis because they shaped the architecture as much
as the process. First, \emph{idempotent seeding} was adopted as a design rule rather
than a convenience: each service ships a deterministic seeder that brings its
database to a known baseline on first boot and that can be re-run safely, so a full
environment---roughly ninety demonstration users, two complete competitions and a
back-filled history of finished matches each with a real starting eleven---can be
reconstructed identically on any host. This made the platform demonstrable on demand
and made defects reproducible against a fixed data set. Second, the project preferred
\emph{reusing existing routes and seams over adding new ones}. When the
database-backed competitions were introduced, for example, they were mounted beneath
an existing gateway route prefix rather than given a new public path, and the
group-chat subsystem was added under the messaging namespace already owned by the
notification service. This conservatism kept the public surface of the gateway small
and stable across iterations, which in turn limited the security and configuration
review that each change required.

\section{System Requirements}
\label{sec:requirements}

Requirements were elicited from the problem analysis and refined during the iterations
described above. They are presented in two parts: functional requirements, organised
by module and by the user roles they serve, and non-functional requirements, which
capture the quality attributes the platform must satisfy.

\subsection{Stakeholders and Roles}
\label{sec:roles}

Before enumerating the requirements it is useful to name the actors the system
serves, because the functional surface is fundamentally role-shaped: what a user may
\emph{see} and what a user may \emph{do} are both derived from the user's role. The
platform recognises a single privileged administrator role and a family of
operational staff and consumer roles. The \textbf{administrator} provisions and
governs all accounts and has all-access. The \textbf{sport manager} owns a sport and
negotiates sponsorship. The \textbf{team manager} administers a specific team's
roster and logistics. The \textbf{coaching staff}---principally the head coach,
supported by assistant and fitness coaches---plan training, set formations and run
matches. The \textbf{medical staff}---doctors and physiotherapists---own the injury
and recovery workflow. \textbf{Scouts} evaluate external prospects and forward
recommendations. \textbf{Sponsors} submit and negotiate contract offers.
\textbf{National-team} records represent external federations that may call up
players. Finally, \textbf{players} and \textbf{fans} are read-mostly consumers who
follow the club, its squad, its competitions and its live matches without the ability
to mutate club data. Each of these roles maps to a realm role in the identity
provider and to a row in the permission map that governs the interface, as described
in Section~\ref{sec:security}.

\subsection{Functional Requirements}
\label{sec:functional-requirements}

MSCMS serves a multi-disciplinary club whose membership spans four football-centric
operational tiers---administration, sport and team management, coaching and medical
staff, and read-only consumers such as players and fans. Access is governed by a
role-based permission map that drives both the navigation a user sees and the
operations a user may perform. The functional requirements are grouped by module; the
modules introduced or substantially reworked in the later iterations---competitions,
live matches, the medical workflow, messaging, notifications, role-based access
control and settings---are highlighted alongside the established modules.

Table~\ref{tab:fr-core} lists the foundational requirements covering identity,
administration, squad management and training. Table~\ref{tab:fr-new} lists the
requirements for the newer modules that distinguish the current system from its
predecessor. The two tables are followed by a narrative elaboration of the
requirements that carry the most behavioural complexity, so that the design choices
made later in the chapter can be traced back to an explicit need.

\begin{table}[htbp]
\centering
\caption{Functional requirements: identity, administration, squad and training modules.}
\label{tab:fr-core}
\renewcommand{\arraystretch}{1.25}
\begin{tabular}{|p{0.07\textwidth}|p{0.20\textwidth}|p{0.60\textwidth}|}
\hline
\textbf{ID} & \textbf{Module} & \textbf{Requirement} \\
\hline
FR-1 & Authentication & The system shall authenticate users against Keycloak using the OAuth~2.0 password grant and issue a JWT access token consumed by the front-end. \\
\hline
FR-2 & Authentication & On successful login the system shall redirect each user to a single dashboard whose content is computed from the user's role. \\
\hline
FR-3 & User management & The administrator shall create, edit, deactivate and delete user accounts across all roles, provisioned through the Keycloak Admin API. \\
\hline
FR-4 & User management & The system shall onboard staff (sport managers, team managers, coaching staff, doctors, physiotherapists, scouts and sponsors) with role assignment. \\
\hline
FR-5 & Club hub & The system shall present a club overview comprising identity, coaching staff, stadium, trophy cabinet, squad and results grouped by competition. \\
\hline
FR-6 & Players \& teams & Authorised staff shall manage sports, teams, rosters, contracts and transfers (incoming, outgoing and national-team call-ups). \\
\hline
FR-7 & Scouting & Scouts shall browse external prospects and author structured scout reports, forwarding recommendations to the sport manager. \\
\hline
FR-8 & Training & Coaching staff shall create training plans, drills and sessions, record attendance, and capture post-session player assessments. \\
\hline
FR-9 & Sponsors & Sponsors shall submit sponsorship offers; the sport manager shall negotiate, accept or reject them with the decision returned to the sponsor. \\
\hline
FR-10 & Reports \& analytics & Analysts shall author match and player analyses and consume the machine-learning predictions for match outcome and player rating. \\
\hline
\end{tabular}
\end{table}

\begin{table}[htbp]
\centering
\caption{Functional requirements: competitions, live matches, medical workflow, messaging, notifications, RBAC and settings.}
\label{tab:fr-new}
\renewcommand{\arraystretch}{1.25}
\begin{tabular}{|p{0.07\textwidth}|p{0.20\textwidth}|p{0.60\textwidth}|}
\hline
\textbf{ID} & \textbf{Module} & \textbf{Requirement} \\
\hline
FR-11 & Competitions & The system shall maintain database-backed competitions, namely La Liga and the new-format UEFA Champions League, each with its teams and fixtures. \\
\hline
FR-12 & Competitions & The Champions League shall be presented as a league phase followed by a knockout bracket, with the knockout stage locked until the league phase is complete. \\
\hline
FR-13 & Competitions & A completed competition (every fixture played) shall become read-only, displaying a ``completed'' badge; ongoing competitions remain editable until the final result is recorded. \\
\hline
FR-14 & Matches & Match creation shall require a starting line-up before the match is persisted, and the scheduled kick-off shall be at least five minutes in the future. \\
\hline
FR-15 & Live match & The system shall drive a live match through a broadcast-style phase machine (first half, half-time, second half, full-time), with extra time and a penalty shoot-out for unresolved knockout ties only. \\
\hline
FR-16 & Live match & During a live match, staff shall record in-game events (goals, cards, substitutions and injuries), with substitutions limited per sport and the line-up locked at creation. \\
\hline
FR-17 & Medical & Medical staff shall manage a Treatment Room that separates currently-injured from recovered players and advances each injured player through a staged journey: injury, diagnosis, treatment, rehabilitation, recovery, fitness test and recovered. \\
\hline
FR-18 & Medical & The system shall maintain general medical catalogues (diagnoses, treatments and recovery programmes) referenced by the per-player journey. \\
\hline
FR-19 & Messaging & The system shall provide real-time team group chat with a general group for all users, user-created groups with invitations (invited/member states), threaded replies and emoji reactions. \\
\hline
FR-20 & Messaging & Chat delivery shall use a WebSocket channel with a polling fallback when the socket is unavailable. \\
\hline
FR-21 & Notifications & The system shall generate event-driven alerts (match scheduled or completed, goal scored, result available, injury, transfer and cancelled training) surfaced through a top-bar bell and a notification centre. \\
\hline
FR-22 & RBAC & The system shall enforce per-role page access so that each role sees only the navigation and operations permitted to it, with players and fans receiving a read-only surface. \\
\hline
FR-23 & Settings & Each user shall manage profile, preference and notification settings with real persistence, and shall change the account password after the current password is verified against Keycloak. \\
\hline
\end{tabular}
\end{table}

\paragraph{Competitions (FR-11--FR-13).} The competitions module is more than a list
of fixtures; it is a small rules engine. When a \emph{league} competition is created,
the back-end pre-generates a complete double round-robin so that every unordered pair
of teams meets the configured number of times with home and away reversed on
alternate rounds. The user therefore enters results only for fixtures that genuinely
exist in the schedule and cannot inflate the standings by replaying the same pairing.
The Champions League is modelled in its current single-league-phase format: a large
group of teams plays a single league stage that is then followed by a two-legged
knockout bracket, and the knockout view is locked until the league phase has been
fully played. Once every fixture of a competition has a recorded result the
competition is considered finished and becomes read-only, surfaced with a
``completed'' badge; deletion by an administrator remains permitted, but result and
date editing are disabled so that historical tables cannot be silently rewritten.

\paragraph{Matches and the live engine (FR-14--FR-16).} A match without a starting
eleven is meaningless, so match creation is an \emph{atomic, line-up-mandatory} flow:
the match, its formation and its starting eleven are persisted together, and the
database is never allowed to hold a match without a line-up. The scheduled kick-off
must be at least five minutes in the future, both to make a deliberate scheduling
decision possible and to give the user time to complete the line-up. At kick-off the
match transitions to live automatically. During the live phase, staff record events
whose recordable minute is bounded by the current phase---a first-half goal cannot be
logged at minute~80, and a regulation half's clock is hard-capped at its length plus
the admin-set stoppage so the displayed minute never runs away into nonsense. The
substitution count is bounded per sport, and recording an injury event flips the
affected player's status, linking the live engine to the medical workflow.

\paragraph{Medical workflow (FR-17--FR-18).} The medical module is organised as a
Treatment Room that partitions players into the currently injured and the recovered,
and it advances each open case through an explicit, staged journey. The
\emph{status of the injury is the single source of truth}: the stage a player is in
is derived from that status rather than tracked independently in the interface, which
guarantees that the Treatment Room and the per-player stepper can never disagree.
General catalogues of diagnoses, treatments and recovery programmes provide reusable
reference data that the per-player journey draws upon.

\paragraph{Messaging and notifications (FR-19--FR-21).} Messaging is a real-time
group-chat subsystem backed by persistent groups: a seeded \emph{general} group to
which every account belongs, plus user-created groups whose invitees remain in an
\emph{invited} state until they accept and become \emph{members}. Delivery favours a
WebSocket push but degrades gracefully to authenticated polling. Notifications and
alerts are event-driven: domain events emitted by the producing services are consumed
asynchronously and turned into alerts that surface through a top-bar bell and a
notification centre with unread badges.

\subsection{Non-Functional Requirements}
\label{sec:nfr}

The non-functional requirements capture the quality attributes that constrain the
design across all modules. They are summarised in Table~\ref{tab:nfr} and elaborated
in the architecture, security and data-design sections that follow. Each
non-functional requirement is paired in the discussion with the concrete mechanism
that satisfies it, so that the quality attributes are demonstrably met by the design
rather than merely asserted.

\begin{table}[htbp]
\centering
\caption{Non-functional requirements grouped by quality attribute.}
\label{tab:nfr}
\renewcommand{\arraystretch}{1.25}
\begin{tabular}{|p{0.07\textwidth}|p{0.21\textwidth}|p{0.59\textwidth}|}
\hline
\textbf{ID} & \textbf{Attribute} & \textbf{Requirement} \\
\hline
NFR-1 & Security & Authentication shall be stateless and JWT-based, with tokens issued by Keycloak and verified at the gateway against the Keycloak JWKS endpoint before any request reaches a domain service. \\
\hline
NFR-2 & Security & Role-based access control shall be enforced in depth: at Keycloak (identity), at the API Gateway (route authorisation) and at the Next.js middleware (navigation guard). \\
\hline
NFR-3 & Scalability & Each domain service shall be independently and horizontally scalable, discoverable through Eureka, with no shared database to create cross-domain contention. \\
\hline
NFR-4 & Scalability & High-volume, cross-domain event fan-out (notifications, alerts) shall be decoupled through asynchronous Kafka messaging rather than synchronous calls. \\
\hline
NFR-5 & Availability & Every service shall expose health endpoints consumable by the orchestrator, and idempotent seeders shall guarantee a consistent baseline state across redeployments. \\
\hline
NFR-6 & Performance & Hot lookups (such as a user profile resolved by Keycloak identifier) shall be cacheable in Redis, and real-time delivery shall favour a WebSocket push over repeated polling. \\
\hline
NFR-7 & Maintainability & Each microservice shall be an independently buildable and releasable module with its own database, Dockerfile and configuration sourced from the Config Server. \\
\hline
NFR-8 & Maintainability & Cross-cutting concerns shall be governed by single sources of truth (the front-end permission map and the back-end seed identifiers) to avoid duplicated, drifting logic. \\
\hline
NFR-9 & Usability & Errors shall be surfaced as friendly, non-technical notifications without exposing stack traces, and the interface shall present a coherent Sportify brand identity. \\
\hline
NFR-10 & Portability & The entire platform shall be reproducible from a single container-orchestration definition so it can be brought up identically on any Docker host. \\
\hline
\end{tabular}
\end{table}

The non-functional requirements are not independent of one another; several reinforce
or tension with their neighbours, and the design resolves those tensions explicitly.
The stateless-token decision (NFR-1) is what makes horizontal scaling cheap (NFR-3),
because a request can be served by any instance of a service without consulting shared
session state. The same decision creates the one genuine exception in the system---the
real-time chat channel, whose handshake cannot carry a bearer header---which is why
that channel authenticates at the message layer and is treated as a deliberately
scoped deviation rather than a hole in the model (Section~\ref{sec:security}).
Asynchronous messaging (NFR-4) trades immediate consistency for decoupling and
resilience: a slow or briefly unavailable notification consumer cannot stall the
user-facing request that produced the event, at the cost of alerts arriving a moment
later. The single-source-of-truth rule (NFR-8) is the maintainability counterpart of
the in-depth authorisation rule (NFR-2): because the same permission map drives both
the rendered navigation and the route guard, the two layers cannot drift apart and
present contradictory access decisions.

\section{System Architecture}
\label{sec:architecture}

MSCMS adopts a \emph{cloud-native microservices} architecture~\cite{microservices}
built on the well-established \emph{API Gateway with Service Discovery and Centralised
Configuration} pattern recommended for Spring Cloud deployments~\cite{springcloud}.
The architecture, illustrated conceptually in Figure~\ref{fig:dashboard} for the
delivered product, separates a single public entry point from a set of privately
networked domain services, each owning its data and lifecycle.

The principal architectural decision was to decompose the system along
operational-domain seams rather than to build a single monolith. Each domain---user
and identity, squad, training and matches, medical, reporting and notifications---has
its own data ownership, release cadence and access patterns. Decomposition allows the
high-traffic services to be scaled independently, allows the schema of one domain to
evolve without coordinating a migration across the entire code base, and allows any
one service container to be rebuilt and redeployed without disturbing the others.
These benefits are realised on a foundation of Spring Boot~3.5.6 on Java~21, in which
each service is an autonomous, embedded-server application~\cite{springboot}.

The seams were chosen along the lines where the domain language and the data change
together. Identity, the squad, the calendar of training and matches, the medical
record, analytical reporting and communication each have a distinct vocabulary, a
distinct set of authoritative tables and a distinct cadence of change, and they are
operated by distinct groups of users. Cutting the system along those lines keeps each
service \emph{cohesive}---its operations all concern one domain---and keeps the
services \emph{loosely coupled}, interacting only through well-defined HTTP contracts
and asynchronous events rather than through a shared schema. The cost of this choice
is that relationships which would be foreign keys in a monolith become opaque
cross-service references (Section~\ref{sec:data-design}); the benefit is that no
single change touches the whole system and no single database is a bottleneck for all
traffic.

Three infrastructure services constitute the control plane:

\begin{itemize}
  \item \textbf{Eureka service registry} (port 8761). Because Docker assigns dynamic
  container addresses, hard-coded service URLs would be brittle. Every business
  service registers with Eureka on boot, and inter-service calls resolve targets by
  logical service identifier rather than by network address.

  \item \textbf{Spring Cloud Config Server} (port 8082). Configuration---database
  URLs, the JWT issuer, broker addresses---is externalised so that per-environment
  values can change without rebuilding container images, satisfying the maintainability
  requirement NFR-7.

  \item \textbf{API Gateway} (port 8080). The gateway is the only externally exposed
  Spring service. It validates JWT tokens, applies global filters such as CORS and
  request logging, and routes by URI prefix to the appropriate domain service. This
  keeps each business service oblivious to authentication concerns and presents a
  single, stable public surface to the front-end.
\end{itemize}

The platform components surrounding the services are Keycloak~26 (the identity
provider hosting the \texttt{mscms} realm), PostgreSQL~14 (one database per service),
Apache Kafka (the asynchronous event backbone), Redis (caching) and RabbitMQ. The
front-end is a Next.js application that communicates exclusively with the gateway,
and two decoupled Python/FastAPI machine-learning services are reached through a
same-origin proxy rather than through the gateway.

\paragraph{Request lifecycle.} A typical authenticated request follows a uniform path
that embodies the architectural separation of concerns:
\begin{enumerate}
  \item The front-end attaches the JWT from browser storage as an
  \texttt{Authorization: Bearer} header and issues the request to the gateway.
  \item The gateway validates the token signature against Keycloak's JWKS endpoint,
  extracts the realm-role claim, and routes the request by URI prefix (for example,
  \texttt{/players/**} to the player-management service).
  \item The selected domain service executes its business logic against its own
  private PostgreSQL database and returns the response through the gateway.
  \item For events that other domains care about---an injury reported, a match
  finished, a goal scored---the originating service publishes a Kafka topic that the
  notification service consumes to fan out alerts and e-mail, decoupling the producer
  from the consumers in line with NFR-4.
\end{enumerate}

\paragraph{Communication topology.} Two distinct planes therefore coexist. On the
\emph{synchronous} plane, the front-end speaks only to the gateway, the gateway speaks
to the domain services, and the domain services occasionally speak to one another
through Eureka-resolved REST calls when an immediate answer is required---for
example, the notification service asking the user service for the roster of the
general chat group. On the \emph{asynchronous} plane, producer services publish domain
events to Kafka and the notification service consumes them, with no producer ever
blocking on a consumer. The front-end touches one further endpoint outside the
gateway: the raw WebSocket of the notification service, used for real-time chat, which
is discussed as a deliberate exception in Section~\ref{sec:security}. The two
machine-learning services sit on their own plane entirely, reached through a
same-origin proxy so that the browser never makes a cross-origin call and the gateway
remains the single Spring entry point.

\section{Microservices Decomposition}
\label{sec:decomposition}

Beyond the three control-plane services, the system comprises six domain microservices.
Each owns a dedicated PostgreSQL database and exposes a cohesive set of operations
behind the gateway. Their responsibilities are summarised in Table~\ref{tab:services}.

\begin{table}[htbp]
\centering
\caption{Domain microservices, their databases and responsibilities.}
\label{tab:services}
\renewcommand{\arraystretch}{1.3}
\begin{tabular}{|p{0.27\textwidth}|p{0.13\textwidth}|p{0.06\textwidth}|p{0.40\textwidth}|}
\hline
\textbf{Service} & \textbf{Database} & \textbf{Port} & \textbf{Responsibility} \\
\hline
user-management & \texttt{mscms\_user} & 8081 & User profiles for every role, Keycloak Admin-API provisioning, authentication endpoints, and national-team records. \\
\hline
player-management & \texttt{mscms\_player} & 8084 & Sports, teams, rosters, contracts, scouting prospects, external teams, transfers and national-team call-ups. \\
\hline
training-match & \texttt{mscms\_training} & 8087 & Training sessions, plans, drills, attendance and assessments; match fixtures, formations, line-ups, in-game events, reviews, statistics, and the database-backed competitions. \\
\hline
medical-fitness & \texttt{mscms\_medical} & 8086 & Injuries, diagnoses, treatments, rehabilitation and recovery programmes, fitness tests and training loads. \\
\hline
reports-analytics & \texttt{mscms\_reports} & 8083 & Cross-domain analytics: match, player, team and training analytics, scout reports and sponsor contract offers. \\
\hline
notification-mail & \texttt{mscms\_notification} & 8085 & In-app notifications and role-targeted alerts, outbound e-mail, the team group-chat subsystem, and consumption of Kafka domain events. \\
\hline
\end{tabular}
\end{table}

Inter-service communication uses two complementary styles. \emph{Synchronous} REST is
used where an immediate, transactional answer is required; external traffic passes
through the gateway, while internal calls resolve their targets through Eureka by
logical service identifier. \emph{Asynchronous} Kafka topics are used for domain
events that fan out to interested parties: producer services such as training-match
and medical-fitness publish events---\texttt{MatchScheduled}, \texttt{MatchResult\-Available},
\texttt{GoalScored}, \texttt{InjuryReported}, \texttt{TransferApproved} and
\texttt{TrainingCancelled}---which the notification service consumes to create alerts
and dispatch e-mail~\cite{kafka}. This separation keeps user-facing requests fast
while still propagating cross-domain effects reliably.

The newer modules were deliberately placed in the services whose data they extend
rather than in new services, minimising operational surface area. The
\emph{competitions} domain and the \emph{live-match engine} live in the training-match
service because they share the matches and line-ups schema. The staged
\emph{medical injury workflow} lives in the medical-fitness service alongside the
injury catalogue. The \emph{team group chat} and the \emph{alert centre} live in the
notification-mail service, which already owned messages and consumed the events that
drive alerts.

This placement decision is an application of the cohesion principle introduced above
and is the reason the public surface stayed small as the system grew. Because the
competitions endpoints were mounted under a path prefix already routed to the
training-match service, no new gateway route was required and no new authorisation
rule had to be reviewed; the same held for the chat-group endpoints under the
messaging namespace. A new service is justified when a new domain language, a new
authoritative data set and a new scaling profile appear together; none of the later
modules met that bar, and adding them to existing services avoided the operational
overhead---another database, another deployment unit, another registration with the
registry---that a new service would have incurred for no architectural gain.

\section{Security and Authentication}
\label{sec:security}

Security is treated as an architectural concern rather than a per-service afterthought.
Identity is centralised in Keycloak, authentication is stateless and token-based, and
authorisation is enforced in depth across three layers.

\subsection{Identity Provider and Token Issuance}

Keycloak hosts the \texttt{mscms} realm, bootstrapped at first boot from a realm
definition that declares the platform's roles, the public OIDC client used by the
front-end, and the seeded demonstration users~\cite{keycloak}. Authentication uses the
OAuth~2.0 \emph{resource-owner password-credentials grant}: the front-end submits the
user's credentials to the gateway, which forwards them to the user-management service;
that service exchanges them at Keycloak's token endpoint and returns the issued JWT
access token~\cite{jwt}. The front-end stores the access token for subsequent calls
and records the user's realm role for the navigation guard.

A refresh token is issued alongside the access token. When an access token expires and
a request consequently fails, the front-end can exchange the refresh token at the
token endpoint for a freshly rotated pair and transparently retry the original
request; only if the refresh token is itself expired or revoked is the user returned
to the login page with storage cleared. This keeps sessions long-lived from the user's
perspective without ever extending the lifetime of the short-lived access token that
the services actually trust.

\subsection{Stateless Verification at the Gateway}

Every subsequent request carries the JWT as an \texttt{Authorization: Bearer} header.
The gateway validates the token's signature against Keycloak's JWKS (JSON Web Key Set)
endpoint and extracts the realm-role claim before routing the request. Because
verification is performed against the published signing keys, no shared session state
is required and the domain services remain stateless, satisfying NFR-1 and supporting
the horizontal-scaling requirement NFR-3. The signing keys are fetched once and cached,
so the common case adds no network round-trip to verification; the gateway re-fetches
only when it encounters a key identifier it has not seen, which is what allows Keycloak
to rotate its keys without a coordinated restart of the platform.

\subsection{Role-Based Access Control in Depth}

Authorisation is enforced through role-based access control (RBAC) at three
layers~\cite{rbac}, providing defence in depth so that the same permission rules apply
whether a user navigates through the interface or calls the API directly:

\begin{enumerate}
  \item \textbf{Keycloak} assigns each user one or more realm roles embedded in the
  JWT.
  \item \textbf{The API Gateway} authorises routes against the role claim, rejecting
  requests whose role is not permitted for the requested prefix.
  \item \textbf{The Next.js middleware} guards navigation client-side using a single
  permission map---the authoritative table of which roles may reach which dashboard
  routes---so that unauthorised pages are never rendered.
\end{enumerate}

The permission map encodes an expert, per-role design: administrators receive
all-access; coaching, medical, scouting, sponsor and analyst roles each receive a
tailored subset; and players and fans receive a read-only surface in which add, edit
and delete controls are suppressed. Because the map is the single source of truth for
both the sidebar that renders the navigation and the middleware that guards it, the
two cannot drift apart (NFR-8).

It is important that the three layers are complementary rather than redundant
substitutes for one another. The middleware guard is a \emph{usability} control: it
prevents a user from navigating to a page they could not use and avoids rendering a
view that would only fail on its first API call. It is not, and is not relied upon as,
a security boundary, because client-side code can always be bypassed. The genuine
security boundary is the gateway, which rejects an unauthorised API call regardless of
how it was issued; even a user who forced their browser past the middleware would
receive a forbidden response from the gateway. Keycloak underpins both by being the
sole authority on who a user is and which roles they hold. The defence-in-depth posture
is therefore that the interface declines gracefully, the gateway refuses
authoritatively, and the identity provider remains the single source of truth for
identity and role.

\subsection{Real-Time Channel and Password Verification}

Two security-sensitive specifics deserve note. First, the team group chat uses a raw
WebSocket endpoint exposed directly by the notification service at
\texttt{ws://\dots/ws/chat}~\cite{websocket}. This path sits deliberately outside the
\texttt{/messages} REST namespace and is reached on the service's own port rather than
through the gateway's JWT filter, because the browser WebSocket API cannot attach a
bearer header during the handshake; the application therefore authenticates the socket
at the message layer and falls back to authenticated polling when the socket is
unavailable.

Second, changing the account password is a two-step, Keycloak-verified operation. The
user-management service first \emph{verifies the current password} by attempting a
password-grant token exchange against Keycloak with the supplied current credentials;
only if that exchange succeeds does the service set the new password through the
Keycloak Admin API's reset-password endpoint. This guarantees that a stolen access
token alone cannot be used to change a password without knowledge of the existing one.
Concretely, the service resolves the user's login identifier (preferring the username
and falling back to the e-mail address), submits it with the supplied current password
to Keycloak's token endpoint, and proceeds only on a successful exchange; a blank or
incorrect current password is rejected with a friendly error before any change is
attempted, and only then is the new password written through the Admin API.

\section{Data Design}
\label{sec:data-design}

The data layer follows a strict \emph{database-per-service} (polyglot persistence)
strategy on PostgreSQL~14~\cite{postgres}. Each of the six domain services owns exactly
one database, and \emph{no physical foreign key crosses a database boundary}. This
isolation is what makes the services independently deployable and independently
scalable, and it directly underpins the scalability and maintainability requirements
NFR-3 and NFR-7.

\subsection{Cross-Service Identifiers}

Because relationships cannot be expressed as cross-database joins, three classes of
identifier are shared as opaque references across the services:

\begin{itemize}
  \item \textbf{Keycloak UUID} (\texttt{keycloakId}) --- the canonical identity of a
  person. Every service stores this string rather than a numeric key into the user
  database, so that a player, coach or sponsor is referenced consistently everywhere.
  \item \textbf{Numeric domain identifiers} --- such as \texttt{teamId},
  \texttt{playerId} and \texttt{matchId}. These are JPA-generated primary keys; the
  seeders preserve a deterministic insertion order so that, for instance, the team
  identifier for FC~Barcelona is stable across every service that references it.
  \item \textbf{Realm role names} --- string identifiers such as \texttt{ADMIN} and
  \texttt{HEAD\_COACH} carried in the JWT and used uniformly for authorisation.
\end{itemize}

The consequence of using opaque references rather than enforced foreign keys is that
referential integrity \emph{across} services becomes an application responsibility
rather than a database guarantee. The design accepts this trade deliberately: the
alternative---a shared schema with real foreign keys---would couple the services back
into a distributed monolith and defeat the independence the decomposition was meant to
achieve. Within a single database, by contrast, ordinary foreign keys are used freely,
so each service still enjoys full relational integrity over the data it owns. The
deterministic seeding rule is the practical safeguard that keeps the cross-service
numeric identifiers aligned, and the use of the Keycloak UUID for people means that
the one identifier most often shared---the identity of a person---is globally unique by
construction and never depends on insertion order at all.

\subsection{Entity-Relationship Model}

Within each database the schema is a conventional relational model managed through
Spring Data JPA and Hibernate. Table~\ref{tab:schema} summarises the principal tables
per database; the complete entity-relationship model is maintained as a dedicated
diagram in the project's documentation.

\begin{table}[htbp]
\centering
\caption{Principal tables per service database.}
\label{tab:schema}
\renewcommand{\arraystretch}{1.3}
\begin{tabular}{|p{0.24\textwidth}|p{0.68\textwidth}|}
\hline
\textbf{Database} & \textbf{Principal tables} \\
\hline
\texttt{mscms\_user} & \texttt{user}, plus a per-role profile extension (player, head\_coach, doctor, physiotherapist, fitness\_coach, team\_manager, sport\_manager, scout, sponsor, fan, national\_team). \\
\hline
\texttt{mscms\_player} & \texttt{sports}, \texttt{teams}, \texttt{rosters}, \texttt{player\_contracts}, \texttt{outer\_teams}, \texttt{outer\_players}, \texttt{player\_transfer\_incoming}, \texttt{player\_transfer\_outgoing}, \texttt{player\_call\_up\_requests}. \\
\hline
\texttt{mscms\_training} & \texttt{training\_sessions}, \texttt{training\_plans}, \texttt{training\_drills}, \texttt{training\_attendance}, \texttt{player\_training\_assessments}, \texttt{matches}, \texttt{match\_formations}, \texttt{match\_lineups}, \texttt{match\_events}, \texttt{player\_match\_statistics}, and the competition tables (\texttt{competition}, \texttt{competition\_team}, \texttt{competition\_fixture}). \\
\hline
\texttt{mscms\_medical} & \texttt{injuries}, \texttt{diagnoses}, \texttt{treatments}, \texttt{rehabilitations}, \texttt{recovery\_programs}, \texttt{fitness\_tests}, \texttt{training\_loads}. \\
\hline
\texttt{mscms\_notification} & \texttt{notifications}, \texttt{alerts}, \texttt{messages}, and the group-chat tables (\texttt{chat\_group}, \texttt{chat\_group\_member}). \\
\hline
\texttt{mscms\_reports} & \texttt{match\_analysis}, \texttt{player\_analytics}, \texttt{team\_analytics}, \texttt{training\_analytics}, \texttt{scout\_reports}, \texttt{sponsor\_contract\_offers}. \\
\hline
\end{tabular}
\end{table}

The schema accommodates the newer modules with minimal disruption. In the user
database, the central \texttt{user} entity is specialised by a one-to-one profile
extension per role, so role-specific attributes (a coach's licence level, a player's
preferred position and market value) live with their role without polluting the common
record. In the training database, the competition tables relate a competition to its
participating teams and to its ordered fixtures, where a fixture's round
distinguishes the Champions League league phase from its knockout rounds. In the
notification database, a chat group owns its members, each carrying an invited or
member status, and the general group is seeded so that every account belongs to it on
arrival.

\subsection{Per-Service Schema Walk-through}
\label{sec:schema-walkthrough}

It is worth describing the principal entities and relationships of each database in a
little more depth, because the shape of the data explains the shape of the operations.

\paragraph{User and identity (\texttt{mscms\_user}).} The common \texttt{user} record
holds the \texttt{keycloakId}, the human-readable identifiers (username, e-mail), basic
demographics and the role. Each role then has a one-to-one extension table that
\emph{extends} the user with its own attributes: a player extension records date of
birth, nationality, preferred position, market value, kit number and the player's
current status, together with the roster and contract identifiers that point into the
player database; a head-coach extension records the sport and team served, the staff
role, years of experience and coaching licence level; a national-team extension records
the federation name, country and contact. This vertical-extension pattern keeps the
common identity in one place while letting each role carry exactly the fields it needs.

\paragraph{Squad (\texttt{mscms\_player}).} A \texttt{sport} owns many \texttt{teams},
and a team owns many \texttt{rosters}, each roster row binding a player (by Keycloak
UUID) to a team for a season. Contracts are held per player by Keycloak UUID with start
and end dates, salary and release clause. The transfer market is modelled with two
symmetric tables: \texttt{player\_transfer\_incoming} links an external player belonging
to an external team to one of the club's teams, while \texttt{player\_transfer\_outgoing}
moves one of the club's players to an external team; both carry a status that walks
through pending, accepted and declined. External teams and their external players are
first-class so that scouting and transfers have concrete targets, and a separate
call-up table records national-team requests for a player.

\paragraph{Training and matches (\texttt{mscms\_training}).} A \texttt{training\_session}
owns its \texttt{training\_drills}, its \texttt{training\_attendance} rows (one per
player, with a present/absent/late/excused/injured status) and its per-player
\texttt{player\_training\_assessments}. Training plans group sessions toward a goal over
a date range. On the match side, a \texttt{match} records the home team, the external
opponent, type, status, sport, venue, competition, season, the score and the kick-off
and finish times. A \texttt{match\_formation} captures the chosen shape and tactical
approach, and it owns the \texttt{match\_lineup} rows that place each player in a
position with a starting, substitute or reserve status and record whether and when they
were substituted. \texttt{match\_event} rows record the in-game timeline---type, minute
and any stoppage---and \texttt{player\_match\_statistics} captures per-player numbers
with sport-specific columns. The competition tables sit alongside: a \texttt{competition}
owns its \texttt{competition\_team} rows and its ordered \texttt{competition\_fixture}
rows, where each fixture carries its two teams, an optional group, a matchday, a round
label and, once played, a score.

\paragraph{Medical (\texttt{mscms\_medical}).} The \texttt{injury} is the aggregate
root: it owns its \texttt{diagnoses}, \texttt{treatments} and \texttt{rehabilitations},
and a rehabilitation in turn owns its \texttt{recovery\_programs}. The injury carries
the status that drives the entire workflow (active, recovering, recovered) together with
the injured body part, severity and dates. Independent of a specific injury, the
database also stores \texttt{fitness\_tests} and \texttt{training\_loads}, which record
objective measurements (test results, session intensity and load, heart-rate figures)
used both in fitness assessment and as input to analytics.

\paragraph{Reporting (\texttt{mscms\_reports}).} This database holds derived and authored
artefacts rather than transactional state. \texttt{match\_analysis}, \texttt{player\_analytics},
\texttt{team\_analytics} and \texttt{training\_analytics} aggregate the operational data
of other domains over a period, several of them storing sport-specific or trend data as
embedded JSON so the schema does not need a column per sport. \texttt{scout\_reports}
attach structured ratings and a sign-or-not recommendation to an external player, and
\texttt{sponsor\_contract\_offers} record a sponsor's offer to a team with an amount,
duration, terms and a negotiation status.

\paragraph{Communication (\texttt{mscms\_notification}).} A \texttt{notification} is a
per-recipient item with a type, category, title, body and read state; an \texttt{alert}
is a higher-level, optionally role-targeted item with a priority, an acknowledgement
state and a resolution state. A \texttt{message} belongs to a chat group, carries its
sender by Keycloak UUID, and may reference a parent message to form a thread. The
\texttt{chat\_group} and \texttt{chat\_group\_member} tables model membership, the latter
carrying the invited-or-member status and who issued the invitation.

\section{Domain and Class Model}
\label{sec:domain-model}

The relational schema described above is materialised in the running services as a set
of JPA entity classes, related by composition and reference. Describing the model at
the class level clarifies which relationships are owning (a parent that contains and
manages the lifecycle of its children) and which are mere references (a foreign-key-like
pointer that the application resolves on demand).

In the \textbf{squad} model, \texttt{Sport} composes many \texttt{Team} objects, and a
\texttt{Team} composes many \texttt{Roster} entries. \texttt{PlayerContract},
\texttt{OuterTeam} and \texttt{OuterPlayer} round out the domain, with an
\texttt{OuterTeam} composing its \texttt{OuterPlayer} squad. The two transfer types,
\texttt{PlayerTransferIncoming} and \texttt{PlayerTransferOutgoing}, reference the teams
and players involved and carry a \texttt{RequestStatus}, expressing the market as a pair
of directed, status-bearing associations rather than as mutable fields on the player.

In the \textbf{match and training} model, \texttt{TrainingSession} composes its
\texttt{TrainingDrill}, \texttt{TrainingAttendance} and assessment children, while
\texttt{Match} composes its \texttt{MatchEvent} and \texttt{PlayerMatchStatistics}
children. Crucially, \texttt{MatchFormation} composes the \texttt{MatchLineup} entries:
the starting eleven is owned by the formation, which is why creating a match, its
formation and its line-up is a single atomic act and why a match cannot exist without a
line-up. \texttt{MatchEvent} carries the enumerated event type and the minute, which the
live engine bounds to the current phase.

In the \textbf{medical} model the composition hierarchy directly encodes the recovery
journey: \texttt{Injury} composes its \texttt{Diagnosis}, \texttt{Treatment} and
\texttt{Rehabilitation} children, and a \texttt{Rehabilitation} composes its
\texttt{RecoveryProgram}. Because the stage of recovery is read from the injury's status
rather than tracked separately, the aggregate root remains the single source of truth and
the staged journey is a projection of that one field. \texttt{FitnessTest} and
\texttt{TrainingLoad} stand apart from the injury hierarchy as independent measurement
entities.

Throughout the model, relationships that would cross a service boundary are expressed not
as object references but as the opaque identifiers described in
Section~\ref{sec:data-design}---a \texttt{playerKeycloakId} string or a numeric
\texttt{teamId}---so the in-memory object graph of any one service never reaches into the
data of another. This is the class-level expression of the database-per-service rule and
is what keeps the services genuinely independent at runtime.

\section{API Surface Overview}
\label{sec:api-surface}

Each domain service exposes a cohesive REST surface behind the gateway, organised by
resource and following conventional HTTP semantics: a collection path supports listing
and creation, an item path supports retrieval, update and deletion, and purpose-built
sub-paths express domain operations that are not pure CRUD~\cite{restapi}.
Table~\ref{tab:api} summarises representative endpoints per service; the prefixes shown
are those routed by the gateway.

\begin{table}[htbp]
\centering
\caption{Representative API endpoints per service (selected, not exhaustive).}
\label{tab:api}
\renewcommand{\arraystretch}{1.25}
\begin{tabular}{|p{0.26\textwidth}|p{0.66\textwidth}|}
\hline
\textbf{Service / prefix} & \textbf{Representative endpoints} \\
\hline
user-management \newline \texttt{/auth}, \texttt{/users}, \texttt{/staff} &
\texttt{POST /auth/login}, \texttt{POST /auth/refresh}, \texttt{POST /auth/admin/create-user}; \texttt{GET/POST /users}, \texttt{GET/PUT/DELETE /users/\{id\}}, \texttt{PUT /users/\{id\}/password} (current-password verified); role collections such as \texttt{/players}, \texttt{/staff}, \texttt{/national-teams}. \\
\hline
player-management \newline \texttt{/players area} &
\texttt{GET/POST /sports}, \texttt{/teams}, \texttt{/rosters}, \texttt{/player-contracts}; \texttt{/outer-teams}, \texttt{/outer-players}; transfer flows \texttt{/transfers/incoming}, \texttt{/transfers/outgoing} and \texttt{/call-ups}, each with status transitions. \\
\hline
training-match \newline \texttt{/training\dots}, \texttt{/matches} &
training: \texttt{/training-sessions}, \texttt{/training-plans}, \texttt{/training-drills}, \texttt{/training-attendance}, \texttt{/player-training-assessments}; matches: \texttt{/matches}, \texttt{/match-formations}, \texttt{/match-lineups}, \texttt{/match-events}, \texttt{/player-match-statistics}; external sync \texttt{POST /matches/sync-external} and standings \texttt{GET /matches/standings/\{code\}}; competitions under \texttt{/matches/leagues} (\texttt{GET/POST}, \texttt{GET /\{id\}}, \texttt{POST /import}, \texttt{POST /\{id\}/fixtures} and \texttt{/fixtures/bulk}, \texttt{PUT /\{id\}/fixtures/\{fid\}}). \\
\hline
medical-fitness \newline \texttt{/injuries} area &
\texttt{GET/POST /injuries}, \texttt{GET/PUT/DELETE /injuries/\{id\}}; the staged children \texttt{/diagnoses}, \texttt{/treatments}, \texttt{/rehabilitations}, \texttt{/recovery-programs}; plus \texttt{/fitness-tests} and \texttt{/training-loads}. \\
\hline
reports-analytics \newline \texttt{/analytics} area &
\texttt{/match-analysis}, \texttt{/player-analytics}, \texttt{/team-analytics}, \texttt{/training-analytics}; \texttt{/scout-reports}; \texttt{/sponsor-contract-offers} with accept/reject transitions. \\
\hline
notification-mail \newline \texttt{/notifications}, \texttt{/alerts}, \texttt{/messages} &
\texttt{/notifications} and \texttt{/alerts} (with \texttt{GET /alerts/target/\{keycloakId\}} and \texttt{GET /alerts/unacknowledged}); messaging \texttt{/messages}; chat groups under \texttt{/messages/groups} (\texttt{GET ?user=}, \texttt{POST}, \texttt{GET /\{id\}/members}, \texttt{GET /invitations}, \texttt{POST /\{id\}/accept}, \texttt{/decline}, \texttt{/invite}); the raw WebSocket \texttt{ws://\dots/ws/chat} sits outside the gateway. \\
\hline
\end{tabular}
\end{table}

Three conventions in this surface are worth highlighting because they recur across
services. First, \emph{literal sub-paths take precedence over identifier patterns}: the
chat-group routes are mounted at the literal \texttt{/messages/groups} so that they are
matched ahead of the \texttt{/messages/\{id\}} item route, which lets a new feature reuse
an existing gateway route without ambiguity. Second, \emph{state transitions are modelled
as explicit endpoints} rather than as opaque field updates---accepting an invitation,
declining it, recording a fixture result, or acknowledging an alert each has its own
verb-bearing path, which makes the API self-describing and keeps authorisation rules
crisp. Third, \emph{external data enters only through the back-end}: the external-sync and
standings endpoints on the training-match service are the sole points at which
third-party data is ingested, after which the front-end reads it from the platform's own
model like any other resource (Section~\ref{sec:tech-stack}).

\section{Technology Stack}
\label{sec:tech-stack}

The platform is assembled from a coherent set of mainstream, well-supported
technologies, summarised in Table~\ref{tab:stack}. The choices reflect the
non-functional priorities: a battle-tested back-end framework with first-class cloud
support, a modern reactive front-end, and standard infrastructure that runs identically
inside containers~\cite{docker}.

\begin{table}[htbp]
\centering
\caption{Technology stack across the back-end, front-end and infrastructure tiers.}
\label{tab:stack}
\renewcommand{\arraystretch}{1.3}
\begin{tabular}{|p{0.18\textwidth}|p{0.32\textwidth}|p{0.40\textwidth}|}
\hline
\textbf{Tier} & \textbf{Technology} & \textbf{Role in MSCMS} \\
\hline
Back-end & Spring Boot 3.5.6 / Java 21~\cite{springboot} & Domain microservices and embedded application servers. \\
\hline
Back-end & Spring Cloud (Gateway, Eureka, Config)~\cite{springcloud} & Routing, service discovery and centralised configuration. \\
\hline
Back-end & Spring Data JPA / Hibernate & Object-relational mapping over PostgreSQL. \\
\hline
Identity & Keycloak 26~\cite{keycloak} & OIDC/OAuth~2.0 identity provider, realm \texttt{mscms}, JWT issuance~\cite{jwt}. \\
\hline
Front-end & Next.js 16 (App Router)~\cite{nextjs} & Server- and client-rendered user interface and route middleware. \\
\hline
Front-end & React 19~\cite{react} & Component model for the dashboard and interactive views. \\
\hline
Front-end & Tailwind CSS v4 & Utility-first styling and the Sportify visual identity. \\
\hline
Real-time & Native WebSocket~\cite{websocket} & Team group chat with a polling fallback. \\
\hline
Data & PostgreSQL 14~\cite{postgres} & One isolated relational database per domain service. \\
\hline
Messaging & Apache Kafka~\cite{kafka} & Asynchronous domain-event backbone for alerts and notifications. \\
\hline
Cache / queue & Redis, RabbitMQ & Hot-lookup caching and transactional messaging. \\
\hline
External data & football-data.org REST API~\cite{footballdata,restapi} & Source of real competition fixtures, results and standings, ingested by the back-end. \\
\hline
Machine learning & Python 3, FastAPI & Decoupled match-outcome and player-rating prediction services. \\
\hline
Platform & Docker \& Docker Compose~\cite{docker} & Reproducible container orchestration of the full stack. \\
\hline
\end{tabular}
\end{table}

A noteworthy data-integration decision is that the back-end is the sole consumer of
the external football data: the training-match service ingests fixtures, results and
standings from the football-data.org REST API on a schedule and persists them into the
matches schema~\cite{footballdata,restapi}, so the front-end never contacts the
external service directly. This keeps the platform's data model authoritative,
shields the user interface from third-party rate limits and outages, and preserves the
single-public-surface property of the gateway architecture. The ingestion is
defensive in practice: it pulls the competition match and standings collections rather
than per-team endpoints (which the free tier rejects), filters to the club's fixtures,
de-duplicates against an external identifier so re-runs do not create duplicate matches,
and reads standings through without storing them so that the live table always reflects
the source.

\section{Behavioural Flows and Sequence Designs}
\label{sec:flows}

The static architecture described above is exercised by a number of behavioural flows.
This section develops the most important of them in detail. Each flow is presented as a
short narrative followed by a numbered sequence of interactions among the participating
components, so that the cooperation between the front-end, the gateway, the domain
services, Keycloak and Kafka is explicit. The flows are illustrated through the
delivered interface in Chapter~\ref{ch:results}.

\subsection{Authentication and Session Establishment}
\label{sec:flow-login}

Authentication establishes a stateless session: the user proves their identity once,
and a signed token thereafter authorises every request without any server-side session
store.

\begin{enumerate}
  \item The user submits a username and password on the Sportify sign-in page.
  \item The front-end issues \texttt{POST /auth/login} to the gateway with the
  credentials.
  \item The gateway forwards the request to the user-management service.
  \item The user-management service performs an OAuth~2.0 password-grant exchange at
  Keycloak's token endpoint (\texttt{/realms/mscms/protocol/openid-connect/token},
  \texttt{grant\_type=password}).
  \item Keycloak validates the credentials and returns an access token, a refresh token
  and an expiry.
  \item The user-management service returns the access token and the user's role to the
  gateway, which relays them to the front-end.
  \item The front-end decodes the token to read the subject (the Keycloak UUID) and the
  realm roles, stores the access token and the Keycloak UUID, and records the role both
  in local storage and in a cookie that the route middleware can read.
  \item The user is redirected to the single \texttt{/dashboard} route, whose content is
  then computed from the stored role.
  \item On every subsequent request the front-end attaches the token as
  \texttt{Authorization: Bearer}; the gateway verifies the signature against Keycloak's
  cached JWKS, checks expiry, extracts the realm-role claim, and routes the request---no
  server-side session is consulted at any point.
\end{enumerate}

When an access token expires, a request fails with an unauthorised status; the
front-end then exchanges the refresh token for a rotated pair and retries the original
request, or, if the refresh token is itself invalid, clears storage and returns the user
to the login page. The result is a session that feels continuous to the user while the
tokens the services actually trust remain short-lived.

\subsection{Scheduling an Official Match from a Competition Fixture}
\label{sec:flow-schedule}

Scheduling an official match is a guarded, atomic flow that begins from a real
competition fixture and ends with a persisted match, formation and starting eleven and a
result written back to the source fixture.

\begin{enumerate}
  \item From a competition view the user selects an unplayed fixture of a still-running
  competition and chooses to schedule it. The interface deep-links into match creation,
  carrying the competition and fixture identifiers and the two teams' names so the
  team-selection steps are pre-filled and skipped.
  \item The user sets a kick-off time. The interface validates that the time parses and is
  at least five minutes in the future, nudging the picker's minimum accordingly, and
  rejects a past or too-soon kick-off with a friendly message.
  \item The user builds the line-up formation-first, choosing a shape and then placing the
  starting eleven by drag-and-drop, with substitutes and reserves designated.
  \item Only once the line-up is complete is the match persisted. The front-end creates the
  match, then its formation, then its line-up rows in one logical act, so the database
  never holds a match without a starting eleven. The fixture's competition and identifier
  are recorded with the match so the result can later be written back.
  \item Because the kick-off is in the future, the match is stored as scheduled and shows a
  live countdown on its card.
  \item At the exact kick-off moment the match transitions to live automatically---the user
  never flips it manually---and the live-match engine takes over its phase machine
  (Section~\ref{sec:flow-live}).
  \item When the match ends, the engine writes the final score back to the originating
  competition fixture, marking it played, so the standings recompute and the fixture is
  shown as completed.
\end{enumerate}

\subsection{Live-Match Phase Progression with Extra Time and Penalties}
\label{sec:flow-live}

The live engine drives a match through a broadcast-style sequence of phases. Control is
gated until kick-off; thereafter an authorised user advances the phases and records
events, while the displayed clock is derived from the kick-off timestamp so the timeline
and the clock always agree.

\begin{enumerate}
  \item At kick-off the match enters \textsc{First Half}. The clock counts from zero, and
  the regulation minute is hard-capped at 45 plus any admin-set stoppage so the display
  reads, for example, ``45+3''' rather than running away into nonsense.
  \item Events recorded during a phase are bounded to that phase's minute range: a
  first-half event must fall in minutes 1--45 (plus stoppage), a second-half event in
  46--90 (plus stoppage). The event is persisted through the match-events endpoint.
  \item The user advances to \textsc{Half Time} (no running clock), then to
  \textsc{Second Half} (counting from minute 45 with its own stoppage), then to
  \textsc{Full Time}.
  \item For an ordinary (league or friendly) match, \textsc{Full Time} ends the match.
  Recording the end posts a result-available alert and writes the score back to the
  source fixture.
  \item For a \emph{knockout} tie that is level at \textsc{Full Time}, the engine continues
  into \textsc{Extra Time}: a first extra period (minutes 91--105), an extra-time
  half-time interval, and a second extra period (minutes 106--120), each a playing phase
  with its own bounded minute range.
  \item If the tie is still level after extra time, the engine enters \textsc{Penalties}, a
  shoot-out with no running clock; the user records the shoot-out outcome.
  \item Throughout, recording a goal posts a real goal-scored alert, and ending the game
  posts a result-available alert and persists the final score. Recording an in-match injury
  event flips the affected player's status to injured, opening a case in the medical
  workflow (Section~\ref{sec:flow-medical}). The substitution count is bounded per sport.
\end{enumerate}

This phase machine is the reason the extra-time and penalty logic is confined to knockout
ties: only a tie that must produce a winner can progress beyond full time, so the engine
checks both the match type and the level scoreline before offering the extra-time
transition. The same minute-bounding logic also lets the system reconstruct the correct
phase for any stored event when a live match is reloaded, by mapping the event's minute
back to the phase whose range contains it.

\subsection{The Injury-to-Recovery Medical Workflow}
\label{sec:flow-medical}

The medical workflow advances an injured player through a staged journey whose stage is
always derived from the injury's status, so the Treatment Room and the per-player stepper
cannot disagree.

\begin{enumerate}
  \item An injury is opened---either logged directly by medical staff or raised as an
  in-match injury event by the live engine. A case is created in the medical-fitness
  service with an active status, and the player appears in the Treatment Room's
  currently-injured list.
  \item Opening the case publishes an injury event to Kafka; the notification service
  consumes it and raises an alert for the relevant staff (Section~\ref{sec:flow-notify}).
  \item Medical staff record a \textsc{Diagnosis} against the injury, drawing on the general
  diagnosis catalogue and adding notes and recommendations.
  \item A \textsc{Treatment} is started, with a type, description and date range; the
  injury's status moves to recovering.
  \item A \textsc{Rehabilitation} plan is created and run by a physiotherapist, owning one or
  more \textsc{Recovery Programs} drawn from the recovery catalogue.
  \item A final \textsc{Fitness Test} confirms readiness, recording an objective result and a
  result category.
  \item On success the injury's status is set to recovered; because the stage is read from
  that status, the player automatically moves from the currently-injured list to the
  recovered list and the stepper shows the journey complete.
\end{enumerate}

The workflow demonstrates the platform's defining characteristics together: a domain-owned
aggregate whose status is the single source of truth, an event published across a service
boundary to drive an alert, and a role-restricted surface in which only medical staff may
advance the stages while everyone else sees the outcome read-only.

\subsection{Real-Time Group Messaging over WebSocket with Polling Fallback}
\label{sec:flow-chat}

Group messaging combines a persistent group model with a real-time delivery channel that
degrades gracefully when the socket is unavailable.

\begin{enumerate}
  \item On opening the messages view, the front-end requests the user's groups
  (\texttt{GET /messages/groups?user=}); the response always includes the seeded general
  group plus every group in which the user is a confirmed member.
  \item For the general group the member roster resolves to every user in the system; the
  service attempts an internal call to the user service and, if that call is rejected, the
  front-end resolves the roster from its own user directory so the general group is never
  empty.
  \item The front-end opens a WebSocket to the notification service at
  \texttt{ws://\dots/ws/chat}. Because the browser WebSocket handshake cannot carry a bearer
  header, the channel is authenticated at the message layer rather than at the gateway.
  \item While the socket is open, sent messages are delivered in real time and incoming
  messages are pushed to all participants. Messages are de-duplicated on arrival (by server
  identifier, then client identifier, then sender-and-content) so a message is never rendered
  twice.
  \item If the socket cannot be established or drops, the client falls back to authenticated
  polling of the messages endpoint, so conversation continues---more slowly---without the
  socket.
  \item A user may create a group, becoming its first member, and invite others; invitees
  remain in an invited state and appear in their pending-invitations list until they accept,
  at which point they become members; declining removes the membership. Replies reference a
  parent message to form threads, and emoji reactions are attached to individual messages.
\end{enumerate}

\subsection{Event-Driven Notification Creation via Kafka}
\label{sec:flow-notify}

Notifications and alerts are produced asynchronously so that a user-facing request is never
delayed by the work of informing other parties.

\begin{enumerate}
  \item A producing service performs a domain action with cross-domain significance---the
  training-match service schedules or completes a match or records a goal, the medical service
  opens an injury, the player service approves a transfer, or a training session is cancelled.
  \item Having committed its own transaction, the producer publishes a corresponding domain
  event (for example \texttt{MatchScheduled}, \texttt{GoalScored},
  \texttt{MatchResultAvailable}, \texttt{InjuryReported}, \texttt{TransferApproved} or
  \texttt{TrainingCancelled}) to a Kafka topic and returns to its caller immediately.
  \item The notification service's event consumer reads the event from Kafka independently of
  the original request.
  \item The consumer creates an alert of the appropriate type and priority, targeting either a
  specific user (by Keycloak UUID) or a role, and persists it; where configured it also queues
  an outbound e-mail.
  \item The front-end's notification centre, opened from the top-bar bell, polls for alerts and
  notifications on a short interval, showing unread badges and letting the user acknowledge an
  alert or mark a notification read.
\end{enumerate}

Because the producer never blocks on the consumer, a slow or briefly unavailable notification
path cannot stall the action that generated the event; the alert simply arrives a moment
later. This is the concrete realisation of the asynchronous-decoupling requirement NFR-4 and is
the mechanism that lets a single live-match goal both update the score and, independently,
inform every interested user.
